\documentclass[10pt,a4paper]{amsart}
\usepackage[margin=19mm]{geometry}
\usepackage{amsmath,amssymb,mathtools}
\usepackage[scr=rsfs,cal=euler]{mathalfa}
\usepackage{xfrac}
\usepackage[unicode,hidelinks,bookmarks=false]{hyperref}
\newcommand{\ie}{i.e.\ }
\newcommand{\cf}{cf.\ }
\newcommand{\resp}{respectively\ }
\newcommand{\Subsection}{\subsection}
\providecommand{\nobreakdash}{\nobreak\hbox{-}}
\setlength{\emergencystretch}{2em}
\multlinegap0pt
\usepackage{bm}
\usepackage{bbm}
\usepackage{dsfont}
\usepackage{xspace}
\usepackage{cancel}
\usepackage{mathtools}
\usepackage{amsthm}
\newtheorem{theorem}{Theorem}
\title[Sobriety of Scott topologies under countability conditions]{Sobriety of Scott topologies under countability conditions}
\author{Zhengmao He, Zhaochen Dong, Kaiyun Wang}
\date{Source: arXiv:2609.24001v2 (CC BY 4.0)}
\begin{document}
\maketitle

\noindent\emph{Source and licence.} Sobriety of Scott topologies under countability conditions. Version arXiv:2609.24001v2 (CC BY 4.0), distributed under the Creative Commons Attribution 4.0 International licence, as recorded in the arXiv licence metadata for this exact version and shown on the abstract page https://arxiv.org/abs/2609.24001v2. Full rights notice: licenses/arxiv-2609.24001-CC-BY-4.0.txt.\par\smallskip

\noindent\textit{Editorial scope.} Counted unit: the Theorem whose source label is \texttt{None} in the article (source0.tex lines 252--257, source label \texttt{None}), delivered together with the article's own complete original proof of it (source0.tex lines 259--336, one \texttt{proof} environment of 78 lines). No mathematics is written, repaired or re-derived: statement and proof are the article's own text line for line. Required context retained uncounted: none. Every definition, display and label of those lines is retained unchanged. Omitted from this package and not redistributed: 1--251, 258--258, 337--492. Nothing inside a retained excerpt is omitted or altered: every retained body is delivered exactly as the article's lines are written, so no omission receipt of any kind accompanies this package. Citations: the retained excerpts carry no citation command at all, so no bibliography entry of the article is transcribed and no reference is redistributed. Reference adaptation: none was needed and none was made; every reference inside the delivered unit resolves inside it, so no unresolved reference or citation is produced. Printed slips kept literal: no printed word, symbol, hypothesis or number of the counted statement or of its proof was altered. Layout and class: the article's own class is \texttt{elsarticle} with packages including enumerate, amssymb, lipsum, geometry, booktabs, amsmath, url, enumitem, graphics, array, dsfont, xy, amsthm, tikz; the wrapper is the lane's pinned \texttt{amsart} editorial wrapper at 10pt on A4, which restates the theorem-like environments the retained text needs and adds the article's own macro definitions that the retained text needs (none), with extra wrapper packages (bm, bbm, dsfont, xspace, cancel, mathtools). The delivered numbering is the unit's own. Assets: no retained span contains a figure environment, a graphics-inclusion command, a PDF-page inclusion, a table or a TikZ picture; no image file is copied into this package, the package declares no asset dependency and no image file is redistributed with it. Licence: version arXiv:2609.24001v2 of this article is distributed under the Creative Commons Attribution 4.0 International licence, as recorded in the arXiv licence metadata for this exact version and shown on the abstract page https://arxiv.org/abs/2609.24001v2. A full rights notice is delivered at licenses/arxiv-2609.24001-CC-BY-4.0.txt. Characters: every retained excerpt is pure ASCII, so the delivered engine, XeLaTeX with its default Latin Modern text font, reports no missing character.
\begin{theorem}
{\rm Let $P$ be a countable dcpo. If
$\Sigma P$ is core-compact, then $P$ is
quasicontinuous. Consequently, $\Sigma P$ is locally compact
sober.}
\end{theorem}

\begin{proof} Fix $x\in U\in\sigma(P)$. By core-compactness, there exists
$V\in\sigma(P)$ such that
$$
x\in V\ll_{\sigma(P)}U,
$$
where $\ll_{\sigma(P)}$ denotes the way-below relation in the
open-set lattice $(\sigma(P),\subseteq)$.

$\mathbf{Claim}$: There exists a finite set $F\subseteq U$ such that $V\subseteq\uparrow F\subseteq U$.

If $U$ is finite, take $F=U$. Otherwise, we enumerate $U=\{u_1,u_2,\ldots\}$.
Assume that no such finite set exists.
Since $U$ is an upper set, for each $n\geq1$ we can choose
$$
a_n\in V\setminus\uparrow\{u_1,\ldots,u_n\}.
$$
For each $k\geq1$, define
$$
C_k=(P\setminus U)\cup\bigcup_{n\geq k}\downarrow a_n.
$$

It is clearly a lower set. Let $D\subseteq C_k$ be directed.
If $\sup D\notin U$, then $\sup D\in C_k$.
Suppose that $\sup D\in U$. Since $U$ is Scott open,
there exists
$$
d_0=u_j\in D\cap U.
$$
The set $D'=D\cap\uparrow d_0$ is a cofinal directed subset
of $D$, and hence
$$
\sup D'=\sup D.
$$
Moreover, $D'\subseteq U$ because $U$ is an upper set.
For every $d\in D'$,  $d\in C_k\cap U$ yields
some $n\geq k$ such that $d\leq a_n$.
Thus $u_j=d_0\leq d\leq a_n$.
The choice of $a_n$ forces $n<j$, so
$$
D'\subseteq\bigcup_{k\leq n<j}\downarrow a_n.
$$
Since $D'$ is nonempty and $\bigcup\limits_{k\leq n<j}\downarrow a_n$ is Scott closed,
, we have
$$
\sup D=\sup D'
\in\bigcup_{k\leq n<j}\downarrow a_n
\subseteq C_k.
$$
This proves that $C_k$ is Scott closed.

Next, we have
$$
\bigcap_{k\geq1}C_k=P\setminus U.
$$
Indeed, $P\setminus U\subseteq\bigcap_{k\geq1}C_k$ is immediate.
For the reverse inclusion, take any $u_j\in U$.
For every $n\geq j$, the choice of $a_n$ gives
$u_j\nleq a_n$. Hence $u_j\notin C_j$.

Put $W_k=P\setminus C_k$.
Then $\{W_k\}_{k\geq1}$ is an increasing sequence of Scott-open
sets satisfying
$$
\bigcup_{k\geq1}W_k=U.
$$
Since $V\ll_{\sigma(P)}U$, there exists $k$ such that
$V\subseteq W_k$. However, $a_k\in V\cap C_k$, a contradiction.
This claim is proved.

We have therefore shown that, for every $x\in U\in\sigma(P)$,
there exist a Scott-open set $V$ and a nonempty finite set
$F\subseteq U$ such that
$$
x\in V\subseteq\uparrow F\subseteq U.
$$

By Lemma 3.1, $P$ is quasicontinuous. Hence, using Lemma 3.2, $\Sigma P$ is locally compact and sober.
\end{proof}

\end{document}
